\documentclass[a4paper, 12pt]{article}
% math symbols
\usepackage{amssymb}
\usepackage{amsmath}
\usepackage{mathrsfs}
\usepackage{mathseries}


\usepackage[margin = 2cm]{geometry}

\tolerance = 1000
\emergencystretch = 0.74cm



\pagestyle{empty}
\parindent = 0mm

\renewcommand{\coursetitle}{DM/ML}
\setcounter{curtask}{1}

\setmathstyle{\includegraphics[scale = 0.05]{../pics/utia-rest.png}}{Приближенные алгоритмы}{1 курс}


\begin{document}

%\teachermode

\begin{definition*}
    Функцию $f\colon 2^{V} \to \mathbb{R}$ будем называть \deftext{субмодулярной}, если для всех $A, B
    \subseteq V$ выполняется:
    $$
        f(A) + f(B) \ge f(A \cup B) + f(A \cap B).
    $$
\end{definition*}

\task{
    Пусть $G \coloneqq (V, E)$~--- неориентированный граф. Для $S \subseteq V$ определим $\delta(S)$ как
    разрез между $S$ и $V \setminus S$, т.е. $\delta(S) = \{ (u, v) \in E \mid u \in S, v \notin S\}$.
    Пусть $c$~--- пропускная способность ребер $G$. Покажите, что функция
    $f(S) \coloneqq \sum\limits_{e \in \delta(S)} c(e)$ субмодулярна.
}

\task{
    Пусть $G \coloneqq (V, E)$~--- неориентированный граф, $c$~--- пропускные способности ребер, $T$~---
    четное множество вершин. Научимся решать следующую задачу (минимальный $T$-нечетный разрез):
    $$
        \min\limits_{S \subseteq V : |S \cap T| \text{ нечетно}} \sum\limits_{e \in \delta(S)} c(e).
    $$
    Также будем говорить, что $S$ является \deftext{$T$-нечетным}, если $|S \cap T|$ нечетно (что вы
    можете сказать про $V \setminus S$?). Пусть $|T| = k$.
    \begin{enumcyr}
        \item\label{item:cutT} Решите следующую задачу за $\bigO{k \cdot \mathrm{Flow}}$:
            $$
                \min\limits_{S \subseteq V : 0 < |S \cap T| < |T|} \sum\limits_{e \in \delta(S)} c(e)
            $$
            (т.е. в обеих частях разреза должна быть по крайней мере одна вершина из $T$).
        \item Пусть $S$~--- из пункта~\ref{item:cutT} и $|S \cap T|$ нечетно. Покажите, что тогда $S$~---
            искомый разрез.
        \item Пусть $S$~--- из пункта~\ref{item:cutT} и $|S \cap T|$ четно. Покажите, что тогда
            существует минимальный $T$-нечетный разрез $A$, такой что $A \subseteq S$ или $A \subseteq
            V\setminus S$. \hinttext{Подсказака:} предлагается рассмотреть произвольный минимальный
            $T$-нечетный разрез и, воспользовавшись субмодулярностью, получить из него искомый разрез
            $A$.
        \item Придумайте, как, воспользовавшись утверждением предыдущего пункта, свести задачу к
            рекурсивному вызову для новых $T$ меньшего размера: приведите алгоритм, работающий за время
            $\bigO{k^2 \cdot \mathrm{Flow}}$.
    \end{enumcyr}
    $\mathrm{Flow}$~--- время работы алгоритма поиска максимального потока.
}
    

\libproblem{algorithms}{boruvka-time}
\libproblem{algorithms}{boruvka-time-planar}
\libproblem{algorithms}{boruvka-improvement}
\libproblem{algorithms}{maj-tree-calc}


\end{document}



%%% Local Variables:
%%% mode: latex
%%% TeX-master: t
%%% End:
